\chapter{Conclusions and Outlook}
\label{cha:conclusions}
    \epigraph{%
    Von guten Mächten wunderbar geborgen,\\
    erwarten wir getrost, was kommen mag.\\
    Gott ist mit uns am Abend und am Morgen\\
    und ganz gewiss an jedem neuen Tag.}
    {Dietrich Bonhoeffer}

%This is left here as an example structure for the conclusions, which should not exceed 3-4 pages in length.

\section{Synopsis}
    \label{sec:concl-synopsins}
    %
    \noindent This thesis has explored the role of the Hubbard manifold in determining the outcome of Hubbard-corrected \ac{DFT} calculations.

%%%%%%%%%%%%%%%%%%%%%%%%%%%%%%%%%%%%%%%%%%%%%%
    %
\section{Emergent themes}
    \label{sec:concl-themes}
    %
    Several recurring themes emerge from this work that merit further reflection.
    %
%%%%%%%%%%%%%%%%%%%%%%%%%%%%%%%%%%%%%%%%%%%%%%
\section[Limitations and perspectives]{Limitations and methodological perspectives}
    \label{sec:concl-limitations}
    %
    The themes discussed above---particularly the continuous nature of localization and the importance of structural complexity---directly inform the limitations of the OR-DFT+U approach.
    %
%%%%%%%%%%%%%%%%%%%%%%%%%%%%%%%%%%%%%%%%%%%%%%
\section{New possibilities}
    \label{sec:concl-possibilities}
    %
    Already today, OR-DFT+U constitutes a valuable addition to the ever-growing zoo of electronic structure methods, enabling access to the longer time and length scales required for \textit{ab initio} molecular dynamics and Monte Carlo approaches.
    %
    \acused{GPU}
    The \ac{GPU}-accelerated diffusion study represents only an initial demonstration---one can now envision realistic simulations of thermally activated phase transitions, or of electrochemical double-layer formation at PBA electrode-electrolyte interfaces.
    %
